\section{Experiments}
\label{sec:experiments}

\subsection{Setup}
\label{sec:setup}

\paragraph{Shapes.}
We evaluate on four \emph{primitive} targets --- \emph{sphere}, \emph{cylinder},
\emph{box}, \emph{square pyramid} --- and two \emph{composite} CSG targets that
test negative-feature machining: a \emph{through-hole sphere} (\texttt{sphere\_hole})
$S_{\mathrm{sph}} \cap \neg S_{\mathrm{cyl}}$ and a \emph{concave bowl}
(\texttt{sphere\_bowl}) $S_{\mathrm{sph}} \cap \neg (S_{\mathrm{sub}} \cap
\{z \le z_c\})$. The composite targets combine a positive outer surface with a
negative interior feature (a through-hole and a hemispherical cavity
respectively), exercising the method's ability to machine concavity and
internal topology from geometry alone. All targets are exposed to the method
only as a pointwise SDF; no result in this section used a shape-dependent
hyperparameter, and the same operating point trains primitives and composites
alike.

\paragraph{Scenario.}
Default stock is a $1~\mathrm{in}^3$ cube at $0.5~\mathrm{mm}$ voxels
($51^3$ grid). Tool radius $r{=}3.175~\mathrm{mm}$, flute height
$h{=}50~\mathrm{mm}$ (Sec.~\ref{sec:gap}, Mechanism 3). Trajectory length
$T{=}128$ waypoints, $N_{\mathrm{iters}}{=}5000$ Adam iterations at
$\alpha{=}10^{-3}$, gradient clip $c_{\mathrm{grad}}{=}0.5$, step
$\Delta t{=}0.45$. Each configuration is run over 3 seeds; we report the
deployable Dice (exact boolean carve, viz stage) mean $\pm$ standard deviation.

\paragraph{Metric.}
All reported numbers are \emph{deployable} Dice: the exact boolean carve
(Eq.~\ref{eq:hard-carve}) of the best-checkpoint trajectory from the canonical
start, after collision-safe truncation. Because $h{=}50~\mathrm{mm}$ clears the
holder throughout the $1~\mathrm{in}$ envelope, viz Dice equals train Dice (no
truncation) for all primitive and composite results. All runs are fully
deterministic given the seed.

\paragraph{Configurations.}
To isolate the contribution of the surrogate-correction mechanisms
(Sec.~\ref{sec:gap}), every results table reports two configurations that differ
\emph{only} in how the soft surrogate is treated. Everything else --- the
simulator, the trajectory parametrization, the optimization procedure, the
evaluation protocol --- is held identical across the two.

\textbf{Shared experimental setup (identical in Base and Full).} Both
configurations use: (i) the same differentiable CSG carving simulator
(Sec.~\ref{sec:simulator}); (ii) the same \emph{random} trajectory
initialization ($\sigma_{\mathrm{init}}{=}0.05$, Sec.~\ref{sec:param}) --- no
shape-dependent init in either; (iii) the same optimization procedure: $5000$
Adam iterations at $\alpha{=}10^{-3}$, $\ell_2$ gradient clip
$c_{\mathrm{grad}}{=}0.5$, step $\Delta t{=}0.45$, $T{=}128$ waypoints; (iv) the
same speed-limit / $z$-floor kinematic gates; (v) the same best-checkpoint
selection over the exact boolean carve (Sec.~\ref{sec:optim}); and (vi) the same
$h{=}50~\mathrm{mm}$ tool-length reachability correction (Mechanism~3). The
deployable evaluation --- hard boolean carve, Dice/ASD/HD95 --- is also identical
and, being $k$-independent, scores both configurations on the same metric. A
third reference point, \emph{Init (iter~0)}, is the deployable Dice of the random
trajectory before any gradient step; because the hard carve is $k$-independent,
this value is shared by both configurations.

\textbf{Methodological differences (Base vs Full).} The two configurations differ
in exactly two places, both acting on the \emph{training-time soft surrogate} and
neither touching the deployable evaluation:
\begin{itemize}\itemsep0pt
  \item \emph{Soft-union sharpness schedule.} Base holds the smoothness $k$ fixed
        at $10$ for the entire run, so the training surrogate has a constant,
        large over-erosion bias ($\mathcal{O}(\log 2 / k)$ per union). Full
        \emph{anneals} $k$ linearly from $k_{\mathrm{init}}{=}10$ to
        $k_{\mathrm{final}}{=}150$ over training (Mechanism~1): a smooth,
        well-conditioned surrogate early for the random init to descend, sharpened
        late so the terminal gradient aligns with the deployable metric.
  \item \emph{Loss de-biasing.} Base evaluates the loss on the soft stock field
        as-is ($\Delta{=}0$). Full adds a constant bias $\Delta{=}-0.5$ voxels to
        the stock SDF before the loss occupancy (Mechanism~2), compensating the
        residual over-erosion so the soft-loss optimum targets the hard carve.
\end{itemize}
No other methodological difference exists between them: the same code path, the
same hyperparameters, the same seeds. The $\Delta$ column in each table is the
per-shape gain of Full over Base and therefore isolates the contribution of these
two surrogate corrections alone. Base is, in effect, the ablation ``what if we
optimze the soft surrogate as-is, with no correction''; Full is the proposed
method. Unless stated otherwise, both are run with identical hyperparameters
across every shape, including the composite CSG targets.

\subsection{Main Result}
\label{sec:main-result}

Table~\ref{tab:main} reports the deployable Dice of the Full method against the
Base configuration on the $1~\mathrm{in}$ stock, with the iteration-0 Dice of the
random trajectory as a reference (Sec.~\ref{sec:setup}, Configurations). Because
Base and Full share every component except the two surrogate corrections, the
$\Delta$ column isolates their effect. The shape-blind method raises mean
deployable Dice from $0.664$ to $0.852$ ($+0.188$), with the largest gains on the
shapes where the fixed-$k$ surrogate most badly misleads the optimizer (sphere
$+0.19$, pyramid $+0.39$). Cylinder and box, which have $z$-invariant
cross-sections and thus a softer structural ceiling, already scored near $0.8$ at
$k{=}10$ and gain correspondingly less. Two observations bear on the
methodological contrast: (i) Base itself moves the init only modestly ($+0.052$
mean, $0.612 \to 0.664$) --- plain soft-surrogate optimization largely treads
water because its gradient is misaligned with the deployable metric; (ii) Full
adds $+0.188$ on top of Base, so the gain is attributable to the surrogate
corrections (annealing $k$ and de-biasing the loss), not to the optimization
machinery, which is already present in Base. The two composite CSG targets
(through-hole sphere, concave bowl) are reported separately in
Sec.~\ref{sec:composite}, as their negative features require distinct geometry.

\begin{table}[t]
  \centering
  \caption{Deployable Dice (3-seed mean $\pm$ std) on the $1~\mathrm{in}$ stock. \emph{Init}: iteration-0 random trajectory (pre-optimization, $k$-independent). \emph{Base}: $k{=}10$, $h{=}50~\mathrm{mm}$, random init. \emph{Full}: + \texttt{k-anneal} ($k_{\mathrm{final}}{=}150$) + \texttt{loss-shift} ($\Delta{=}-0.5$).}
  \label{tab:main}
  \small
  \begin{tabular}{lccc c}
    \toprule
    Shape & Init (iter 0) & Base ($k{=}10$) & Full method & $\Delta$ \\
    \midrule
    Sphere   & 0.550 & 0.638 & $\mathbf{0.826 \pm 0.018}$ & $+0.188$ \\
    Pyramid  & 0.386 & 0.427 & $\mathbf{0.813 \pm 0.003}$ & $+0.386$ \\
    Cylinder & 0.707 & 0.774 & $\mathbf{0.905 \pm 0.010}$ & $+0.131$ \\
    Box      & 0.803 & 0.816 & $\mathbf{0.865 \pm 0.009}$ & $+0.049$ \\
    \midrule
    Mean     & 0.612 & 0.664 & $\mathbf{0.852}$           & $+0.188$ \\
    \bottomrule
  \end{tabular}
\end{table}

\subsection{Ablation: Closing the Soft/Hard Gap}
\label{sec:ablation}

Table~\ref{tab:ablation} isolates the three mechanisms of
Sec.~\ref{sec:gap}. Starting from the original short-tool fixed-$k$ configuration
($h{=}25~\mathrm{mm}$, $k{=}10$, deployable Dice computed after collision
truncation), each mechanism is added cumulatively (seed 1, for comparability).

\begin{table}[t]
  \centering
  \caption{Cumulative ablation, $1~\mathrm{in}$ stock, seed 1. Each row adds one mechanism. Deployable Dice (viz $=$ train when untruncated).}
  \label{tab:ablation}
  \small
  \begin{tabular}{lccccc}
    \toprule
    Configuration & Sph. & Pym. & Cyl. & Box & Mean \\
    \midrule
    $k{=}10$, $h{=}25$ (orig.\ baseline, truncated) & 0.579 & 0.425 & 0.717 & 0.816 & 0.634 \\
    $+ h{=}50$ (reachability)              & 0.638 & 0.427 & 0.774 & 0.816 & 0.664 \\
    $+ k$-anneal $k_{\mathrm{final}}{=}150$ & 0.830 & 0.796 & 0.895 & 0.843 & 0.841 \\
    $+ \Delta{=}-0.5$ (loss-shift)         & 0.840 & 0.817 & 0.895 & 0.865 & 0.854 \\
    \bottomrule
  \end{tabular}
\end{table}

The reachability correction ($h{=}25 \to 50$) removes the truncation ceiling but,
in isolation, barely moves Dice: at $k{=}10$ the optimizer is still pulled toward
the over-eroding soft optimum. Smoothness annealing is the dominant lever
($+0.18$ mean), aligning the terminal gradient with the hard metric. The
de-biasing loss shift adds a further $+0.013$ mean by compensating the residual
over-erosion. The two surrogate corrections are complementary: $k$-anneal
\emph{reduces} the per-union bias; loss-shift \emph{cancels} what remains.

\subsection{Initialization-Mode Ablation}
\label{sec:init-ablation}

Table~\ref{tab:init} confirms the operating-point choice of random
initialization (Sec.~\ref{sec:param}). Without $k$-anneal, the structured
\texttt{raster\_fine} init helps over random (it pre-covers the part envelope,
giving the optimizer a better starting basin when the surrogate itself is
misleading). With $k$-anneal, the gap collapses: the de-biased gradient discovers
the geometry from a random start just as well, so the structured prior becomes
redundant for the primitives. We therefore adopt random as the shape-blind
default. The surface-following \texttt{raster\_fine} init is retained only for
the concave-bowl composite CSG target, where it lifts Dice by $\sim 0.03$ over
random at equal waypoint budget.

\begin{table}[t]
  \centering
  \caption{Initialization mode, $1~\mathrm{in}$ sphere, seed 1. Deployable Dice. With $k$-anneal active, random matches the structured init.}
  \label{tab:init}
  \small
  \begin{tabular}{lcc}
    \toprule
    Init mode & $k{=}10$ (no anneal) & $k$-anneal $k_{\mathrm{final}}{=}80$ \\
    \midrule
    Random (\texttt{random})         & 0.642 & 0.784 \\
    Structured (\texttt{raster\_fine}) & 0.635 & 0.788 \\
    \bottomrule
  \end{tabular}
\end{table}


\subsection{Loss-Shift Sign and Magnitude}
\label{sec:lossshift-ablation}

Table~\ref{tab:lossshift} confirms the de-biasing direction and sweet spot on the
through-hole sphere (feasible-shell geometry; the narrow negative feature makes
the over-erosion bias most visible). Positive $\Delta$ increases over-carve and
lowers Dice; negative $\Delta$ reduces over-carve and lifts Dice, with
$\Delta{=}-0.5$ optimal and $\Delta{=}-1.0$ over-correcting into under-carve.
The sign and order of magnitude are consistent with the predicted
$\Delta \approx -\log 2 \cdot k_{\mathrm{ref}} / k_{\mathrm{final}} \approx -0.24$.

\begin{table}[t]
  \centering
  \caption{Loss-shift sweep on the through-hole sphere (seed 1). Deployable viz Dice and over-carved voxel count (carved minus target).}
  \label{tab:lossshift}
  \small
  \begin{tabular}{lcc}
    \toprule
    $\Delta$ (vox) & Deployable Dice & Over-carve (vox) \\
    \midrule
    $+1.0$  & 0.215 & $+70{,}332$ \\
    $+0.5$  & 0.232 & $+63{,}116$ \\
    $+0.24$ & 0.243 & $+58{,}785$ \\
    $\phantom{+}0$   & 0.237 & --- \\
    $-0.24$ & 0.248 & $+54{,}509$ \\
    $-0.5$  & \textbf{0.246} & $+51{,}375$ \\
    $-1.0$  & 0.232 & $+47{,}152$ (under-carve) \\
    \bottomrule
  \end{tabular}
\end{table}

\subsection{Composite Shapes}
\label{sec:composite}

The two composite CSG targets are reported in Table~\ref{tab:composite}. They are
the strictest test of shape-blindness: their geometry is a non-trivial boolean
combination, yet the method receives only the combined SDF field, with no
indication that a hole or cavity is present.

\begin{table}[t]
  \centering
  \caption{Composite CSG targets, $1~\mathrm{in}$ stock. Deployable Dice (viz $=$ train, no truncation). \emph{Init}: iteration-0 random trajectory (pre-optimization, $k$-independent). \emph{Base}: $k{=}10$, $h{=}50~\mathrm{mm}$. \emph{Full}: + \texttt{k-anneal} + \texttt{loss-shift} $\Delta{=}-0.5$.}
  \label{tab:composite}
  \small
  \begin{tabular}{lccc c l}
    \toprule
    Target & Init (iter 0) & Base ($k{=}10$) & Full method & $\Delta$ & Note \\
    \midrule
    Concave bowl (\texttt{sphere\_bowl}), $T{=}128$ & 0.415 & 0.443 & $\mathbf{0.634 \pm 0.030}$ & $+0.191$ & 3 seeds; ls neutral \\
    Concave bowl (\texttt{sphere\_bowl}), $T{=}256$ & 0.451 & 0.451 & $\mathbf{0.659 \pm 0.001}$ & $+0.208$ & 3 seeds; denser path \\
    Through-hole sphere (\texttt{sphere\_hole})$^\dagger$ & 0.401 & --- & $0.246$ & --- & open frontier \\
    \bottomrule
  \end{tabular}\\
  \small $^\dagger$feasible-shell geometry; see text.
\end{table}

\paragraph{Concave bowl (\texttt{sphere\_bowl}).}
The bowl is machinable: its negative feature is a wide interior cavity that the
swept cylinder can reach, so the loss can penalize the cavity floor and walls.
The full method raises deployable Dice from $0.443$ to $0.634 \pm 0.030$ over 3
seeds ($+0.191$), and a denser $T{=}256$ trajectory reaches $0.659 \pm 0.001$.
The gain over the $k{=}10$ base ($+0.19$) is on the same order as the primitive
sphere, confirming $k$-anneal transfers to composite positive-plus-concave
geometry. Loss-shift is neutral on the bowl: its concavity is wide enough that
soft-union over-erosion does not fill it, so the de-biasing correction has
nothing to cancel --- consistent with the loss-shift analysis of
Sec.~\ref{sec:lossshift-ablation}, where the shift helps only where over-erosion
is the binding constraint.

\paragraph{Through-hole sphere (\texttt{sphere\_hole}).}
The through-hole sphere is the open frontier. With the default geometry the
through-hole shell ($1.9~\mathrm{mm}$) is thinner than the tool
($3.175~\mathrm{mm}$), physically infeasible; we therefore evaluate on a
feasible-shell variant (sub-radius $9.525~\mathrm{mm}$). Even so, the soft
optimum trains to $\approx 0.72$ Dice but deploys at only $\approx 0.25$:
at any finite $k$, soft-union over-erosion fills the narrow negative feature, so
the loss cannot penalize the hole interior --- the optimizer has no gradient
signal for the hole wall. Strikingly, optimization \emph{degrades} the
deployable Dice from the iteration-0 value of $0.401$ to $0.246$: the random
init carves a tentative hole, but following the over-erosion-biased gradient
fills it back in. Loss-shift helps marginally ($0.237 \to 0.246$) by reducing
over-carve around the hole, but the hole interior remains structurally
unpenalized. This is the limit of pure soft-loss methods: a topology-aware or
hard-carve-in-the-loop loss term would be required to close it, which we leave
to future work. The result is included to delineate where the shape-blind
soft-surrogate approach succeeds (positive surfaces, wide concavities) and
where it does not (narrow negative topology).

\subsection{Scaling to Larger Stock}
\label{sec:scaling}

To confirm the method scales beyond the $1~\mathrm{in}$ scenario, we run the full
operating point on a $2~\mathrm{in}^3$ stock with a proportionally scaled
$h{=}100~\mathrm{mm}$ tool (same tool-to-stock reachability ratio) and a
$22.86~\mathrm{mm}$-radius target. Table~\ref{tab:scaling} reports the
deployable Dice (viz $=$ train, no truncation) over seeds. The $k$-anneal win
transfers: sphere at $k{=}10$ stalls at $\approx 0.55$, whereas the full method
reaches $0.758 \pm 0.012$ --- a $+0.22$ gain that matches the $1~\mathrm{in}$
pattern, indicating the method is not tuned to a single absolute scale. The
residual gap to the $1~\mathrm{in}$ sphere ($0.826$) is attributable to the
larger tool-to-stock ratio at $2~\mathrm{in}$, not to a method failure. Box
scales the best of the four: its $z$-invariant cross-section gains from the
finer relative resolution of the larger stock, reaching $0.927$ and exceeding
the $1~\mathrm{in}$ box ($0.865$). Pyramid is the hardest at $2~\mathrm{in}$
($0.577 \pm 0.002$): more sloped surface at the same tool-to-stock ratio leaves
less reachable margin.

\begin{table}[t]
  \centering
  \caption{Scaling to $2~\mathrm{in}^3$ stock, $h{=}100~\mathrm{mm}$ tool (same tool-to-stock ratio). Deployable Dice (viz $=$ train, no truncation), seed mean $\pm$ std.}
  \label{tab:scaling}
  \small
  \begin{tabular}{lcc}
    \toprule
    Shape ($2~\mathrm{in}$) & Base ($k{=}10$) & Full method \\
    \midrule
    Sphere   & $\approx 0.55$ & $0.758 \pm 0.012$ (3 seeds) \\
    Box      & ---            & $0.927$ (2 seeds) \\
    Pyramid  & ---            & $0.577 \pm 0.002$ (3 seeds) \\
    \bottomrule
  \end{tabular}
\end{table}

\subsection{Negative Results and Dead Levers}
\label{sec:negative}

For completeness, Table~\ref{tab:negative} records the configurations that did
\emph{not} improve on the operating point, so that they are not re-explored.

\begin{table}[t]
  \centering
  \caption{Negative results. None improves mean deployable Dice; each is shape-dependent or saturated.}
  \label{tab:negative}
  \small
  \begin{tabular}{lc p{0.5\linewidth}}
    \toprule
    Configuration & Mean Dice & Finding \\
    \midrule
    \textbf{Full method ($T{=}128$)} & \textbf{0.852} & operating point \\
    $T{=}256$ waypoints             & 0.826 & shape-dependent: helps pyramid ($+0.04$), hurts sphere/cyl/box; not a blind win \\
    $k_{\mathrm{final}}{=}500$      & $<0.841$ & high-$k$ hurts positive features; does not close narrow-hole gap \\
    $k_{\mathrm{final}}{=}1000$     & $<0.841$ & gap persists/widens; $k$-anneal cannot close narrow-hole gap alone \\
    $h{=}75~\mathrm{mm}$            & $\approx 0.841$ & $= h{=}50$ within seed noise; tool length saturated \\
    $\Delta > 0$ (loss-shift)       & $<0.237^*$ & wrong direction; monotonically increases over-carve \\
    Air-cut regularizers            & --- & loss-based air reduction is shape-dependent: dice-NEUTRAL $-58\%$ air on sphere ($w_{\mathrm{air}}{=}1$), structural tradeoff on pyramid (Sec.~\ref{sec:trajqual}); free $-45\%$ via checkpoint selection \\
    \bottomrule
  \end{tabular}\\
  \small $^*$on the through-hole sphere.
\end{table}

\subsection{Deployable Trajectory-Quality Measures}
\label{sec:trajqual}

Beyond geometric accuracy, a deployable toolpath must also be \emph{efficient}
and \emph{safe}: it should minimize non-cutting (air) time and avoid tool
breakage. We report three deployable measures computed on the exact boolean
carve (independent of the soft surrogate used for optimization): total toolpath
time $t_{\mathrm{tot}}$ (s), air-cutting time $t_{\mathrm{air}}$ (s), and a
trajectory-level tool-breakage probability $P_{\mathrm{brk}}{\in}[0,1]$ from a
per-segment engagement model
$P_{\mathrm{brk}}{=}1{-}\exp(-\sum_t \sigma((\ln \mu_F(t) - \ln F_{\mathrm{ref}})/\sigma_r))$
with $\mu_F(t){=}k_c\, e(t)\, v_{\mathrm{vox}}^3 / (\Delta t\, D)$. These are
\emph{new and uncalibrated}; we treat them as reporting axes alongside Dice.

\paragraph{Soft vs.\ hard Dice.}
The differentiable surrogate optimizes a \emph{soft} (smooth-union) carve; the
deployable outcome is the \emph{hard} (boolean) carve. We report both: soft Dice
is the training metric (the proven $0.85$-class operating point), hard Dice is
the sharp deployable score. At the shape-blind operating point (random init,
fixed $k{=}10$), soft Dice reaches $0.84$--$0.94$ across primitives while hard
Dice is quantized to 1-voxel steps and caps lower ($0.40$--$0.81$); closing the
soft/hard gap is what the $k$-anneal mechanism of Sec.~\ref{sec:gap} achieves.

\paragraph{Air time is bimodal across shapes.}
Table~\ref{tab:trajqual} reports the three measures at the operating point.
The air fraction $t_{\mathrm{air}}/t_{\mathrm{tot}}$ is strongly shape-dependent:
the prismatic shapes (box, cylinder)---whose $z$-invariant cross-sections the
raster init covers with the tool permanently engaged---produce fast, air-free
trajectories ($\le 1\%$ air, $3$--$11~\mathrm{s}$). The freeform surfaces
(sphere, pyramid) require 3D repositioning and incur $21$--$50\%$ air
($20$--$37~\mathrm{s}$). The prior ``$\sim 30\%$ air is inherent'' figure is an
average over this bimodal distribution; the headroom for air-time reduction is
concentrated entirely in the freeform shapes.

\begin{table}[t]
  \centering
  \caption{Deployable trajectory-quality at the shape-blind operating point (seed 1, $1~\mathrm{in}$ stock). Soft Dice is the training metric; hard Dice is the sharp boolean carve. Air fraction is $t_{\mathrm{air}}/t_{\mathrm{tot}}$. Breakage model calibrated with $F_{\mathrm{ref}}{=}5$ (see text).}
  \label{tab:trajqual}
  \small
  \begin{tabular}{lcccc c}
    \toprule
    Shape & Soft Dice & Hard Dice & $t_{\mathrm{air}}$ (s) & $t_{\mathrm{tot}}$ (s) & Air \% \\
    \midrule
    Box      & 0.917 & 0.814 & 0.0  & 3.1  & 0\% \\
    Cylinder & 0.944 & 0.726 & 0.1  & 11.4 & 1\% \\
    Sphere   & 0.842 & 0.565 & 4.1  & 19.6 & 21\% \\
    Pyramid  & 0.892 & 0.397 & 18.3 & 36.5 & 50\% \\
    \bottomrule
  \end{tabular}
\end{table}

\paragraph{Breakage calibration.}
At the default reference force $F_{\mathrm{ref}}{=}50~\mathrm{N}$ the per-step
engagement at $0.5~\mathrm{mm}$ voxels with a $3.175~\mathrm{mm}$ tool is tiny,
so $F_{\mathrm{cut,max}}{=}0$ and $P_{\mathrm{brk}}$ collapses to numerical zero
across all shapes---the model does not fire. Lowering the reference threshold
tenfold ($F_{\mathrm{ref}}{=}5$, $F_{\mathrm{max}}{=}10$) makes it fire:
$P_{\mathrm{brk}}$ rises to $0.004$--$0.008$ and $F_{\mathrm{cut,max}}$ becomes
nonzero, while Dice is unchanged (the breakage loss term is disabled, so
calibration affects only the report). Even calibrated, $P_{\mathrm{brk}}<1\%$
and $F_{\mathrm{cut,max}} \ll F_{\mathrm{max}}$ with $\mathrm{broken}{=}0$
everywhere: at this operating point tool breakage is \emph{not} a binding
constraint, so $P_{\mathrm{brk}}$ is a reporting axis rather than an
optimization target.

\paragraph{Air-time loss: dice-neutral on the sphere, structural on the pyramid.}
The differentiable air-time loss $w_{\mathrm{air}} \cdot \overline{\Delta t \cdot \mathrm{airfrac}}$
(Sec.~\ref{sec:method}) penalizes non-cutting segments. Sweeping its weight
$w_{\mathrm{air}}{\in}\{10^{-3},10^{-2},10^{-1},1,3\}$ on the two freeform shapes
reveals a sharply shape-dependent result. On the \emph{sphere} the loss is a
dice-\emph{neutral} air win: $w_{\mathrm{air}}{=}1.0$ holds Dice at
$0.844 \pm 0.001$ over 3 seeds (vs.\ $0.842$ baseline) while cutting air time
$4.13 \to 1.72~\mathrm{s}$ ($-58\%$) and total time $19.6 \to 13.4~\mathrm{s}$
($-32\%$); $w_{\mathrm{air}}{=}3.0$ pushes air to $0.52~\mathrm{s}$ ($-87\%$)
still dice-neutral. The sphere's $21\%$ air is unnecessary repositioning the loss
removes for free. On the \emph{pyramid} the same loss trades Dice for air:
$w_{\mathrm{air}}{=}1.0$ gives $-39\%$ air but costs $-0.076$ Dice;
$w_{\mathrm{air}}{=}0.3$ gives $-15\%$ air for $-0.027$ Dice. The pyramid's $50\%$
air is \emph{structural}---repositioning between steep slopes the tool must
perform to cover the part. The default $w_{\mathrm{air}}{=}10^{-3}$ is too weak
to move air on any shape; effective reduction needs $w_{\mathrm{air}} \ge 1.0$.
This refines the flat ``$\sim 30\%$ air inherent'' air-regularizer tradeoff of
Sec.~\ref{sec:negative} into a bimodal, shape-dependent answer: air is
removable-for-free on the sphere, structural on the pyramid, and absent on the
prismatic shapes.

\paragraph{Composite checkpoint selection: a free air win.}
A dice-neutral air reduction is also achievable without re-optimizing, through
the composite best-checkpoint criterion
$\mathrm{best\_score} = \mathrm{Dice} - w^{\mathrm{best}}_{\mathrm{air}}\,\tilde t_{\mathrm{air}}$.
With $w^{\mathrm{best}}_{\mathrm{air}}{=}0.05$ and the air-time loss
\emph{disabled} (identical optimization to baseline), selection alone picks a
sphere checkpoint at equal Dice ($0.842$) but $-45\%$ air ($4.13 \to 2.25~\mathrm{s}$)---a
lower-air checkpoint exists in the baseline trajectory that pure-Dice selection
misses. Combining $w_{\mathrm{air}}{=}1.0$ with $w^{\mathrm{best}}_{\mathrm{air}}{=}0.05$
reaches $-78\%$ air dice-neutral. (On the pyramid, selection removes only
$\sim 3\%$ air---structural air is not selectable away.) This makes
checkpoint-selection cost the preferred lever when the baseline trajectory
already contains a low-air checkpoint.
